% SPDX-License-Identifier: Apache-2.0
\section{The Lowering, Checked}
\label{sec:x-checked}

A lowering that is only printed is a claim. From this section on, a
lowered unit is also run, and against the same stimulus the Rust
simulation had, twice. The build does it: the example is run with
\code{TXHDL\_VHDL} and \code{TXHDL\_VERILOG} set, so beside its
trace it writes the unit's VHDL and its Verilog, which
\code{\#[lower]} renders from the same statements, and a ports file;
\code{//tools/fst2tb} reads the trace and the ports and writes a
testbench, in VHDL and again in Verilog, that replays the trace's
inputs into the entity and asserts, tick by tick, that the entity's
registers and outputs are what the trace holds; \code{vhdl\_test}
from \code{rules\_nvc} simulates the VHDL under \code{nvc}, and a
\code{cc\_test} over \code{rules\_verilator} simulates the Verilog
under Verilator, both built from source, hermetically. Each test
passes only if every comparison held, so both emitters are checked,
and a defect in either shows as a disagreement in one of the two.

Both testbenches hold the trace as data rather than as code. A cycle
is a frame: the inputs for its edge, then each value expected after
it with a bit that says whether to check it. The VHDL testbench keeps
the frames as hex literals in packages; the Verilog one reads them
from a file beside it with \code{\$readmemh}, one vector a cycle, and
one loop applies and checks a vector. A testbench is then the same
size however long the run is. Written out a statement per cycle, the
Verilog of the SD host's run was two million lines, which took
Verilator 11~GiB and twelve minutes to compile and was stopped on the
smaller CI runner (issue 601); as a loop it takes 25~MiB and under a
second. A failed comparison names the signal, the tick, the value
expected and the value found.

The testbench's timing is the model of time, made literal. One tick is
one nanosecond and the clock rises at even ticks. An input the trace
holds at tick $2k$ was set before that edge, so it is applied at
$2k-1$. A register the trace holds at $2k$ took its value at that
edge, so it is checked at $2k+1$, through a VHDL-2008 external name.
A wire the Rust process drove at $2k$ was computed from the registers
as they stood before that edge, and from the inputs of that edge, so
it is checked at $2k-1$ after those inputs are applied and a few
deltas have passed, one per level of wires, where the entity's continuous assignment shows the same
value. That last rule
is a fact about the runtime worth knowing: a wire a process drives
inside its iteration lags the hardware's continuous assignment by the
process's own step.

The counter of the cheat sheet, the blinky, the sequencer, the ALU,
the scratchpad, the decoder and the two-process sampler pass.
A negative control was run once: the counter's adder changed to add
two makes the test fail at the first tick, so a pass means something.

The stage of Section~\ref{sec:x-stage} passes too, and that closes a
finding. The runtime's channel was a mailbox: an offer sat until
taken, and the consumer took it in the same step it was offered, so
\code{valid} was never high at a tick's end and \code{ready} meant
the slot was empty, while the lowering's channel was a valid/ready
handshake held across edges; a trace of one was not a check of the
other. The runtime's channel is now an elastic buffer of two with the
handshake registered on both sides, as the runtime document states,
and the netlist's port is what that buffer shows: the head and the
room are inputs the testbench applies a step later than a wire, since
they are registered state, and the take and the offer are outputs
checked as any wire is, the offer's data only under its valid. The
change moved the FIFO's, the join's and the arbiter's traces by the
buffer's latency and changed nothing else about them.

The Verilog check earned its place on its first run. Every unit's
VHDL had agreed with its trace, and the same units' Verilog, from the
same statements, did not, twice, in the RISC-V core alone. An unsized
number in a branch of a conditional inside a concatenation is thirty-two
bits wide in Verilog, so every field above it moved; and signedness in
Verilog is decided by the whole expression, so an arithmetic shift
inside an unsigned conditional chain shifted logically, and one
\code{srai} came out wrong. Neither is a defect VHDL can have, and
neither showed until the Verilog was simulated. Both are fixed in the
emitter, a number sized by the other branch and the shift
self-determined inside \code{\$unsigned}, and a wire the core had
named \code{byte}, a keyword in SystemVerilog, was renamed.
